\author{}
\documentclass[11pt,a4paper]{article}
\usepackage[utf8]{inputenc}
\usepackage[T1]{fontenc}
\usepackage[french]{babel}

% Paquets pour les mathématiques
\usepackage{amsmath, amssymb, array, physics}

% Marges
\usepackage{geometry}
\geometry{hmargin=2cm, vmargin=2cm}

% Paquet pour les boîtes colorées
\usepackage[many]{tcolorbox}

% Définition d'un style personnalisé pour les définitions (bleu)
\newtcolorbox{boiteDef}[1]{
    colback=blue!5!white,
    colframe=blue!75!black,
    title={\textbf{#1}},
    arc=3mm,
    boxrule=1pt,
    fonttitle=\bfseries
}

% Définition d'un style pour les propriétés (rose)
\newtcolorbox{boiteProp}[1]{
    colback=pink!5!white,
    colframe=purple!70!pink,
    title={\textbf{#1}},
    arc=3mm,
    boxrule=1pt,
}

% Définition d'un style pour les rappels (vert)
\newtcolorbox{boiteRappel}[1]{
    colback=green!10!white,
    colframe=green!50!black,
    title={\textbf{#1}},
    arc=3mm,
    boxrule=1pt,
}

% Définition d'un style pour les remarques (jaune)
\newtcolorbox{boiteRemarque}[1]{
    colback=yellow!10!white,
    colframe=yellow!70!black,
    title={\textbf{#1}},
    arc=3mm,
    boxrule=1pt,
}

\renewcommand{\thesection}{\Roman{section}}

\begin{document}

\title{\textbf{Chapitre 7 : Composées, dérivation, continuité}}
\date{}
\maketitle

\section{Composées de fonctions}

\subsection{Définition}

\begin{boiteDef}{Définition :}
    Soit $u$ une fonction définie sur un intervalle $I$ et $v$ une fonction définie sur un intervalle $J$ tel que pour tout $x \in I$, $u(x) \in J$. La fonction composée de $u$ suivie de $v$ est la fonction définie pour tout $x \in I$ par :
    \newline
    $(v \circ u)(x) = v(u(x))$
\end{boiteDef}

\begin{boiteRemarque}{Remarque :}
    \begin{itemize}
        \item On lit "v rond u"
        \item $v \circ u$ est en général différent de $u \circ v$
        \item On peut composer plusieurs fonctions à la suite : $v \circ u \circ w$ etc.
        \item On a déjà utilisé les composées dans les limites sans utiliser le symbole $\circ$.
    \end{itemize}
\end{boiteRemarque}

\subsection{Dérivation et composées}

\begin{boiteProp}{Propriété :}
    Soit $u$ une fonction définie sur un intervalle $I$ et $v$ une fonction dérivable sur un intervalle $J$ tel que pour tout $x \in I$, $u(x) \in J$. 
    \newline
    Alors la fonction $v \circ u$ est dérivable sur $I$ et pour tout $x \in I$ :
    \newline
    $(v \circ u)'(x) = v'(u(x)) \times u'(x)$
\end{boiteProp}

\begin{boiteRappel}{Rappel :}
    Soit $g$ une fonction dérivable sur un intervalle $I$ de $\mathbb{R}$ et soient $a$ et $b$ deux réels. Alors la fonction $f:x \mapsto g(ax + b)$ est dérivable sur $I$ et pour tout $x \in I$ : $f'(x) = a \times g'(ax + b)$
\end{boiteRappel}

\begin{boiteProp}{Conséquences}
    Soit $u$ une fonction dérivable sur $I$ et $n$ un entier relatif non nul.
    \begin{itemize}
         \item La fonction $u^n$ (avec $u(x) \neq 0$ si $n < 0$) est dérivable sur $I$ et $(u^n)'= n \times u' \times u^{n-1}$
         \item La fonction $\sqrt{u}$ (avec $u(x) > 0$) est dérivable sur $I$ et $(\sqrt{u})' = \frac{u'}{2 \sqrt{u}}$
         \item La fonction $e^u$ est dérivable sur $I$ et $(e^u)' = u' \times e^{u}$
    \end{itemize}
\end{boiteProp}

\section{Convexité}

\subsection{Convexité et lecture graphique}

\begin{boiteDef}{Définition :}
    $f$ est une fonction définie sur un intervalle $I$. 
    \newline
    $C$ est sa courbe représentative sur $I$ dans un repère.
    \newline
    \begin{itemize}
        \item $f$ est convexe sur $I$ si pour tous points $A$ et $B$ (distincts) de $C$, le segment $[AB]$ est au-dessus de la courbe $C$ entre $A$ et $B$.
        \item $f$ est concave sur $I$ si pour tous points $A$ et $B$ (distincts) de $C$, le segment $[AB]$ est en dessous de la courbe $C$ entre $A$ et $B$.
    \end{itemize}
\end{boiteDef}

\begin{boiteRemarque}{Remarques :}
    \begin{itemize}
        \item La fonction carrée et la fonction valeur absolue sont convexes sur $\mathbb{R}$.
        \item La fonction racine carrée est concave sur $\mathbb{R}^+$.
    \end{itemize}
\end{boiteRemarque}

\begin{boiteProp}{Propriété :}
    Soit $f$ une fonction dérivable sur $I$ et $C$ sa courbe représentative dans un repère.
    \begin{itemize}
        \item $f$ est convexe sur $I$ si et seulement si la courbe $C$, sur $I$, est entièrement au-dessus de chacune de ses tangentes.
        \item $f$ est concave sur $I$ si et seulement si la courbe $C$, sur $I$, est entièrement en dessous de chacune de ses tangentes.
    \end{itemize}
\end{boiteProp}

\begin{boiteRemarque}{Remarques :}
    \begin{itemize}
        \item La fonction carrée et la fonction exponentielle sont convexes sur $\mathbb{R}$.
        \item La fonction racine carrée est concave sur $\mathbb{R}^+$.
    \end{itemize}
\end{boiteRemarque}

\subsection{Point d'inflexion}

\begin{boiteDef}{Définition :}
    $f$ est une fonction dérivable sur un intervalle $I$ et $C$ sa courbe représentative dans un repère ; $a \in I$.
    \newline
    Le point $A(a, f(a))$ est un point d'inflexion de la courbe $C$ si en $A$, la courbe $C$ traverse sa tangente.
\end{boiteDef}

\begin{boiteRemarque}{Conséquence :}
    En $a$, la fonction $f$ passe de concave à convexe ou inversement.    
\end{boiteRemarque}

\subsection{Convexité et dérivées}

\begin{boiteProp}{Propriété :}
    Soit $f$ une fonction dérivable sur un intervalle $I$.
    \begin{itemize}
        \item $f$ est convexe sur $I$ si et seulement si $f'$ est croissante sur $I$.
        \item $f$ est concave sur $I$ si et seulement si $f'$ est décroissante sur $I$.
    \end{itemize}
\end{boiteProp}

\begin{boiteDef}{Définition :}
    Soit $f$ une fonction dérivable sur un intervalle $I$.
    \newline
    On dit que $f$ est deux fois dérivable sur $I$ si sa dérivée $f'$ est elle-même dérivable sur $I$. 
    \newline
    La fonction dérivée de $f'$, notée $f''$, est appelée fonction dérivée seconde de $f$.
\end{boiteDef}

\begin{boiteProp}{Propriété :}
    Soit $f$ une fonction deux fois dérivable sur un intervalle $I$.
    \begin{itemize}
        \item $f$ est convexe sur $I$ si pour tout $x \in I$, $f''(x) \geq 0$.
        \item $f$ est concave sur $I$ si pour tout $x \in I$, $f''(x) \leq 0$.
    \end{itemize}
\end{boiteProp}

\begin{boiteProp}{Propriété :}
    Si $f''$ est positive alors la courbe représentative de $f$ est au dessus de ses tangentes.    
\end{boiteProp}

\begin{boiteRemarque}{Remarque :}
    On peut démontrer des inégalités en utilisant la convexité d'une fonction.
    \newline
    On retrouve ainsi que $e^x \geq 1 + x$ pour tout $x \in \mathbb{R}$.
    \newline
    En effet, au point d'abcisse $0$, la courbe de la fonction exponentielle admet une tangente d'équation $y = 1 + x$. 
    \newline
    $y = f'(0)(x - 0) + f(0) \equiv y = e^0 \times x + e^0 \equiv y = x + 1$
    \newline
    $f(x) = e^x$ $f'(x) = e^x$ $f''(x) = e^x > 0$ 
    \newline
    La fonction exponentielle vérifie bien $f''(x) > 0$ donc sa courbe est au dessus de ses tangentes.
    \newline
    $e^x \geq 1 + x$.    
\end{boiteRemarque}

\begin{boiteProp}{Propriété :}
    $f$ est une fonction deux fois dérivable sur $I$ et $C$ est sa courbe représentative dans un repère ; $a \in I$.
    \newline
    Le point $A(a, f(a))$ est un point d'inflexion de $C$ si et seulement si $f''$ s'annule en changeant de signe en $a$.    
\end{boiteProp}

\section{Continuité et théorème des valeurs intermédiaires}

\subsection{Continuité d'une fonction}

\begin{boiteDef}{Définition :}
    Soit $f$ une fonction définie sur $I$ et $a \in I$.
    \newline
    \begin{itemize}
        \item $f$ est continue en $a$ si $f$ a une limite en $a$ égale à $f(a)$, c'est-à-dire si $\lim_{x \to a} f(x) = f(a)$
        \item $f$ est continue sur $I$ si $f$ est continue en tout réel $a$ de $I$.
    \end{itemize}
\end{boiteDef}

\begin{boiteRemarque}{Vocabulaire :}
    Si $f$ n'est pas continue en $a$, on dit que $f$ est discontinue en $a$.
\end{boiteRemarque}

\begin{boiteProp}{Propriétés :}
    \begin{itemize}
        \item Les fonctions affines, la fonction carrée, la fonction cube, les fonctions polynômes sont continues sur $\mathbb{R}$.
        \item La fonction inverse est continue sur $]-\infty;0[$ et $]0;+\infty[$.
        \item La fonction racine carrée est continue sur $[0;+\infty[$.
        \item La fonction exponentielle est continue sur $\mathbb{R}$.
        \item Toute fonction définie sur $I$ et obtenue par opérations ou composition des fonctions précédentes est continue sur $I$.
        \item Les fonctions rationnelles sont continues sur tout intervalle contenu dans leur ensemble de définition.
    \end{itemize}
\end{boiteProp}

\begin{boiteProp}{Propriété :}
    Soit $f$ une fonction définie sur $I$ et $a \in I$.
    \begin{itemize}
        \item Si $f$ est dérivable en $a$, alors $f$ est continue en $a$.
        \item Si $f$ est dérivable sur $I$, alors $f$ est continue sur $I$.
    \end{itemize}
\end{boiteProp}

\begin{boiteRemarque}{Remarque :}
    La réciproque de la propriété précédente est fausse : une fonction peut être continue en $a$ sans être dérivable en $a$. 
\end{boiteRemarque}

\subsection{Théorème des valeurs intermédiaires}

\begin{boiteRappel}{Convention :}
    Une flèche dans un tableau de variations d'une fonction $f$ indique :
    \begin{itemize}
        \item la stricte croissance ou stricte décroissance de $f$ sur l'intervalle correspondant
        \item la continuité de la fonction $f$ sur cet intervalle
    \end{itemize}    
\end{boiteRappel}

\begin{boiteProp}{Théorème des valeurs intermédiaires :}
    Soit une fonction $f$ définie et dérivable sur un intervalle $I$ et soient $a$ et $b$ deux éléments de $I$ avec $a < b$.
    \newline
    Pour tout nombre réel $k$ compris entre $f(a)$ et $f(b)$, l'équation $f(x) = k$ admet au moins une solution sur l'intervalle $I$.
    \newline
    Autrement dit, pour tout réel $k$ compris entre $f(a)$ et $f(b)$, il existe au moins un réel $c$ tel que $f(c) = k$.
\end{boiteProp}

\begin{boiteRemarque}{Remarques :}
    Graphiquement, pour tout réel $k$ compris entre $f(a)$ et $f(b)$, la droite d'équation $y = k$ coupe au moins une fois la courbe représentative de $f$.
    \newline
    Le TVI est un théorème d'existence : il affirme l'existence d'au moins une solution d'une équation $f(x) = k$ mais il ne permet pas de connaître le nombre exact de solutions (ni leur valeur exacte).
    \newline
    On va obtenir des valeurs approchées de ces solutions (si on ne peut pas résoudre algébriquement ces équations).
\end{boiteRemarque}

\begin{boiteProp}{Corollaire du théorème des valeurs intermédiaires :}
    Si $f$ est continue et strictement monotone sur un intervalle $[a; b]$, alors pour tout réel $k$ compris entre $f(a)$ et $f(b)$, l'équation $f(x) = k$ admet une unique solution dans l'intervalle $[a; b]$.
\end{boiteProp}

\begin{boiteDef}{Conséquence :}
    Si $f$ est continue et strictement monotone sur $[a; b]$ et si $f(a) \times f(b) < 0$, alors l'équation $f(x) = 0$ admet une unique solution sur $[a; b]$.
\end{boiteDef}

\begin{boiteRemarque}{Remarque :}
    On parle d'intervalle fermé $[a; b]$ mais on peut avoir une (ou les deux) bornes infinies.
\end{boiteRemarque}

\section{Fonctions continues et suites convergentes}

\subsection{Image d'une suite par une fonction}

\begin{boiteDef}{Définition :}
    Soit $(u_n)$ une suite définie sur $\mathbb{N}$.
    \newline
    Soit $f$ une fonction définie sur un intervalle $I$ contenat tous les termes $u_n$ de la suite.
    \newline
    L'image de la suite $(u_n)$ par la fonction $f$ est la suite $(f(u_n))$.
\end{boiteDef}

\subsection{Image d'une suite convergente par une fonction continue}

\begin{boiteProp}{Propriété :}
    Soit $(u_n)$ une suite qui converge vers un nombre réel $l$.
    \newline
    $I$ est un intervalle tel que $l \in I$ et pour tout entier naturel $n$, $u_n \in I$.
    \newline
    Pour toute fonction $f$ définie sur l'intervalle $I$ et continue en $l$, la suite $(f(u_n))$ converge vers le réel $f(l)$.
\end{boiteProp}

\subsection{Suites du type : $u_{n+1} = f(u_n)$}

\begin{flushleft}
    f est une fonction définie et continue sur un intervalle $I$ tel que pour tout réel $x \in I$, $f(x) \in I$.
    \newline
    $(u_n)$ est la suite définie par la donnée de $u_0 \in I$ et pour tout entier naturel $n$, $u_{n+1} = f(u_n)$.
    \newline
    Lorsque la suite $(u_n)$ converge vers un réel $l$, la suite $f(u_n)$ converge vers le réel $f(l)$.
    \newline
    On a : $\displaystyle \lim_{n \to +\infty} f(u_n) = f(l)$, c'est à dire : $\displaystyle \lim_{n \to +\infty} u_{n+1} = f(l)$
    \newline
    Or $\displaystyle \lim_{n \to +\infty} u_{n} = l$, donc $\displaystyle \lim_{n \to +\infty} u_{n+1} = l$.
    \newline
    L'unicité de la limitepermet donc d'affirmer que $f(l) = l$.
    \newline
    On peut conclure que si la suite $(u_n)$ converge vers un réel $l$ de $I$ et si $f$ est continue sur $I$ (donc en particulier en $l$), alors $l = f(l)$.
    \newline
    $l$ est donc solution dans $I$ de l'équation $f(x) = x$.
\end{flushleft}












\end{document}
